\subsection{XLA CPU Workflow: Adapter and Measurement Protocol}
\label{app:xla-workflow}
The workflow experiment uses XLA revision
\code{39c348196c834a7ee45a57f5b44688d53ae99a64}.
Its opt-in hook replaces execution of \code{TransposeFolding} at the pass's
existing CPU-pipeline position, after simplification and before layout
assignment. The baseline invokes the original pass with its original CPU
callbacks. The integrated path first applies the two TEPL transpose/dot rules,
then invokes the original pass for residual unsupported cases. Both use a synchronous PjRt client with one CPU device and the default
CPU lowering options, then execute the remaining CPU optimization pipeline, code generation, and executable
loading. The original pass is unchanged when the hook is unset.

\paragraph{Division of responsibility.}
A shared adapter exposes a bounded HLO graph as dot and transpose nodes with
opaque tensor inputs. The compiled patterns determine the required root
operation kinds and import depth; shared HLO nodes retain their identities.
TEPL's matcher selects captures, checks the written guards and operation
contracts, evaluates the certified attribute program, and builds replacement
metadata. The adapter contains no separate dot/transpose pattern recognizer
or axis-remapping formula. XLA's original CPU callback must approve each
changed dot operand. XLA APIs install the resulting graph, preserving output
shape, layout, precision and derived metadata. Each selected rule applies at
most once to an original root, following the upstream pass's single traversal.

The two rules retain the original nonidentity, rank, single-contraction and
fixed-batch conditions and infer all four RHS dot axis fields per rule.
The bridge supports dense FP32 dot and transpose nodes. Other operations
remain opaque captures; unsupported dtypes, sparse dots and extra-operand
forms remain with XLA. Backend rejections are counted, and errors after an
admitted TEPL match fail the experiment. The residual native fold count is
reported separately in the artifact. This prototype runs one compilation at
a time; it does not assess concurrent compiler requests.

\paragraph{Workloads and inputs.}
We use the public Flax v0.12.0 \code{lm1b} TransformerLM implementation
with one layer, two heads, embedding and QKV dimensions 32, feed-forward
width 64, vocabulary size 128, and sequence length 16. Batches are one and
four, dropout is zero, and evaluation is deterministic. JAX 0.7.2 and Flax
0.12.0 export the initial HLO with parameters as runtime arguments.
The DLRM FP32 training-step and ResNet-50 FP32 prediction HLOs are copied
unchanged from the pinned XLA CPU benchmarks. The workload inventory is
fixed before measuring pass applications; zero-application workloads remain
included. The artifact retains source URLs, exact input bytes and hashes,
model settings, binaries, implementation source, and the hook patch in
\code{evaluation/xla-workflow/}.

These are untrained workload instances. FP32 arguments are small positive
values determined by parameter number, coordinates and seed; integer inputs
are one and Boolean inputs are false. Three seeds vary the floating-point
arguments. This avoids invalid negative variance inputs without claiming
model accuracy or realistic data distributions.

\paragraph{Checks and timing.}
The native bridge is first checked over ranks two through six, both operand
sides and all axis swaps, including zero and singleton extents and unequal
free/contracting extents: 420 transport/exclusion checks and four rejection
controls. Two HLO controls fold both operands, including a shared operand, and compare
the original, upstream and TEPL graphs numerically on three inputs each.
For each model, diagnostic warmups verify and retain HLO before and after
the target pass and after full optimization. Canonical HLO comparison ignores
instruction names, identifiers and debug metadata while retaining constants, backend
configuration and control dependencies. This is structural agreement with
the baseline, not a proof of arbitrary whole-program equivalence.
The baseline and integrated CPU executables are compared on three
input sets, with absolute and relative tolerances $10^{-5}$; FP32 outputs
must be finite.

One untimed warmup per variant precedes ten paired compilation rounds with
alternating order on one Intel Core Ultra 7 155H machine. The integrated
timer includes a fresh portable-program load and validation for every
compilation. Parsing, cloning, HLO comparison and numerical execution are
outside the timer. Diagnostic HLO serialization and extra verification inside
the hook run only during warmup. No match or attribute-output cache is used.
Application counts must be stable across rounds, and TEPL applications plus
residual native folds must equal baseline folds. One application counts one
folded operand, including when upstream folds both operands in one replacement.
The main table gives median paired compilation change and its min--max range;
the following table gives absolute costs. No execution-time speedup is measured.
\begin{table}[!htbp]
\caption{CPU compilation costs in milliseconds: median [min,max] over ten measured rounds per variant. TEPL program loading is included in integrated compilation.}
\label{tab:xla-workflow-timings}
\small
\begin{tabularx}{\linewidth}{@{}Yrrr@{}}\toprule
Workload & Baseline & Integrated & Program load\\\midrule
TransformerLM (1) & \shortstack[r]{86.9\\[-2pt]{\scriptsize[77.6,95.6]}} & \shortstack[r]{104.0\\[-2pt]{\scriptsize[96.4,115.2]}} & \shortstack[r]{15.1\\[-2pt]{\scriptsize[14.1,16.5]}}\\
TransformerLM (4) & \shortstack[r]{115.7\\[-2pt]{\scriptsize[106.8,131.9]}} & \shortstack[r]{132.1\\[-2pt]{\scriptsize[118.3,139.7]}} & \shortstack[r]{16.9\\[-2pt]{\scriptsize[15.5,19.4]}}\\
DLRM training (6,272) & \shortstack[r]{288.1\\[-2pt]{\scriptsize[258.7,327.3]}} & \shortstack[r]{305.9\\[-2pt]{\scriptsize[266.4,345.2]}} & \shortstack[r]{18.5\\[-2pt]{\scriptsize[16.1,23.6]}}\\
ResNet-50 (1) & \shortstack[r]{247.2\\[-2pt]{\scriptsize[235.4,263.4]}} & \shortstack[r]{264.4\\[-2pt]{\scriptsize[246.4,303.2]}} & \shortstack[r]{19.3\\[-2pt]{\scriptsize[18.1,25.3]}}\\
\bottomrule\end{tabularx}
\end{table}

